\documentclass[11pt,twocolumn]{article}

\usepackage[utf8]{inputenc}
\usepackage[T1]{fontenc}
\usepackage{times}
\usepackage{geometry}
\geometry{letterpaper, margin=0.75in, columnsep=0.3in}
\usepackage{graphicx}
\usepackage{booktabs}
\usepackage{amsmath}
\usepackage{hyperref}
\usepackage{xcolor}
\usepackage{listings}
\usepackage{caption}
\usepackage{float}
\usepackage{enumitem}
\usepackage{multirow}
\usepackage{array}

\hypersetup{
    colorlinks=true,
    linkcolor=blue!60!black,
    citecolor=blue!60!black,
    urlcolor=blue!60!black
}

\lstset{
    basicstyle=\ttfamily\scriptsize,
    breaklines=true,
    frame=single,
    backgroundcolor=\color{gray!8},
    keywordstyle=\color{blue!70!black},
    stringstyle=\color{red!60!black},
    commentstyle=\color{green!50!black},
    columns=flexible,
    aboveskip=6pt,
    belowskip=6pt
}

\setlength{\parskip}{4pt}
\setlength{\parindent}{0pt}

\title{\Large\textbf{Structured RPC vs.\ Text-Based Sub-Agent Orchestration:\\A Comparative Analysis of Inter-Agent\\Communication Paradigms}}

\author{
    \textbf{Ope Olatunji}\textsuperscript{1} \quad \textbf{Fola (AI Research Assistant)}\textsuperscript{1}\\[4pt]
    \textsuperscript{1}AgenticMail / OpenClaw Research\\[2pt]
    \texttt{ope.olatunji@outlook.com}
}

\date{March 27, 2026 (Revised)}

\begin{document}
\maketitle

\begin{abstract}
As multi-agent AI systems evolve from single-model architectures to orchestrated ensembles of specialized agents, the communication layer between agents becomes a critical design decision. This paper presents a comparative analysis of two inter-agent communication paradigms: (1)~structured Remote Procedure Call (RPC) task delegation, exemplified by the \texttt{call\_agent} pattern, and (2)~free-text sub-agent spawning, exemplified by the \texttt{sessions\_spawn} pattern. Through extensive empirical testing across \textbf{16 sub-agent tasks spanning 10 industries}---including multi-agent parallel collaboration scenarios---we evaluate both paradigms on latency, output quality, domain versatility, and collaboration capability. We find that structured RPC achieves dramatically lower latency ($\sim$1\,s vs.\ 24--107\,s) and produces machine-parseable outputs, while text-based spawning excels at complex reasoning, creative collaboration, and domain-specialist tasks where output richness matters more than speed. We argue that the optimal architecture is a \textbf{hybrid}: RPC for data retrieval and composable subtasks, spawning for autonomous specialist work and multi-agent collaboration.
\end{abstract}

\textbf{Keywords:} multi-agent systems, inter-agent communication, RPC, orchestration, AI agents, tool use, collaboration

\section{Introduction}

The rapid advancement of large language model (LLM) capabilities has given rise to agentic AI systems---autonomous agents equipped with tools, memory, and the ability to take actions in the world~\cite{yao2023react, wang2024survey}. A natural extension of single-agent systems is multi-agent orchestration, where a primary agent delegates subtasks to specialized subordinate agents. This pattern mirrors microservice architectures in traditional software engineering, where the communication protocol between services is as important as the services themselves.

Despite growing interest in multi-agent frameworks such as AutoGen~\cite{wu2023autogen}, CrewAI, and LangGraph, relatively little attention has been paid to the \emph{communication layer} between agents. Most implementations default to natural language as the inter-agent medium. Our initial study (v1, March~26, 2026) challenged this assumption using a single weather-retrieval task. This revised paper dramatically expands the empirical foundation with \textbf{16 sub-agent tasks across 10 industries}, including \textbf{multi-agent parallel collaboration} scenarios, providing a more nuanced and honest assessment of both paradigms.

We identify five critical axes of comparison:

\begin{enumerate}[itemsep=2pt,topsep=4pt]
    \item \textbf{Latency and throughput}---How quickly can the orchestrator receive actionable results?
    \item \textbf{Output structure and composability}---Can results be programmatically consumed?
    \item \textbf{Orchestrator agency}---Does the pattern preserve or diminish the primary agent's control?
    \item \textbf{Domain depth and reasoning quality}---How well does each paradigm handle complex, specialist tasks?
    \item \textbf{Multi-agent collaboration}---Can agents work on complementary parts of a larger problem?
\end{enumerate}

\section{Background and Related Work}

\subsection{Multi-Agent Architectures}

Multi-agent systems in AI have a long history~\cite{wooldridge2009}, but the LLM era has introduced a new paradigm where agents are language models augmented with tool access. Recent frameworks take varying approaches:

\textbf{AutoGen}~\cite{wu2023autogen} uses conversational message passing in a group chat pattern. \textbf{CrewAI} defines agents with roles and goals, communicating through natural language task delegation. \textbf{LangGraph} models agent interactions as state machines with typed state, offering more structure but still relying on text-based message passing.

None of these frameworks make a strong distinction between structured data exchange and conversational exchange between agents.

\subsection{RPC in Distributed Systems}

Remote Procedure Call (RPC) is a well-established paradigm in distributed computing~\cite{birrell1984}. Modern implementations like gRPC use Protocol Buffers for structured, typed message exchange with strict schemas. The application of RPC principles to LLM-based multi-agent systems remains largely unexplored in the literature.

\section{Methodology}

\subsection{Experimental Setup}

We conducted comparative tests within the OpenClaw agent framework (v2026.3), an open-source personal AI assistant platform. The system runs on consumer hardware (Apple Mac mini, ARM64, macOS~25.2.0, Node.js~v25.5.0) with Claude (Anthropic) as the underlying language model.

\textbf{Method~A: Structured RPC (\texttt{call\_agent}).} The orchestrating agent submits a task with a structured payload to a named subordinate agent, receiving a structured JSON result inline.

\textbf{Method~B: Text-Based Spawning (\texttt{sessions\_spawn}).} Creates an isolated sub-agent session with a natural language prompt. The orchestrator yields its turn, waits for push-based completion, and receives free-form natural language output.

\subsection{Test Battery}

Unlike our initial study which used a single weather-retrieval task, this revision employs a comprehensive \textbf{16-task battery} across 10 domains, organized in three waves executed with up to 5--6 concurrent sub-agents.

\textbf{Wave~1---Single-Agent Domain Specialists} (6 parallel):
\begin{enumerate}[itemsep=1pt,topsep=2pt]
    \item \emph{Finance:} AAPL stock analysis with live web data
    \item \emph{Healthcare:} Drug interaction analysis (Metformin/Lisinopril/Atorvastatin)
    \item \emph{Legal:} SaaS Terms of Service contract risk analysis
    \item \emph{Software Engineering:} Security-focused code review with corrected code
    \item \emph{Marketing:} Multi-channel Gen~Z product launch campaign (\$50K budget)
    \item \emph{Education:} 4-week ML curriculum for working professionals
\end{enumerate}

\textbf{Wave~2---Specialists + Collaboration Seeds} (5 parallel):
\begin{enumerate}[itemsep=1pt,topsep=2pt,start=7]
    \item \emph{Supply Chain:} Shenzhen$\to$Charlotte shipping route optimization
    \item \emph{Data Science:} Sales forecasting model comparison (ARIMA/XGBoost/Prophet/LSTM)
    \item \emph{Academic Research:} RAG literature review with citations
    \item \emph{Creative Writing:} Sci-fi story Part~A (collaborative, sequential handoff)
    \item \emph{Finance (M\&A):} Financial due diligence for \$400M SaaS acquisition
\end{enumerate}

\textbf{Wave~3---Multi-Agent Collaboration} (5 parallel):
\begin{enumerate}[itemsep=1pt,topsep=2pt,start=12]
    \item \emph{Creative Writing:} Sci-fi story Part~B (continues Part~A output)
    \item \emph{Legal (M\&A):} Legal due diligence (complements \#11)
    \item \emph{Operations (M\&A):} Operational due diligence (complements \#11 + \#13)
    \item \emph{Biostatistics:} Phase~III clinical trial statistical analysis plan
    \item \emph{Regulatory/Clinical:} Clinical trial regulatory strategy (complements \#15)
\end{enumerate}

\subsection{Evaluation Criteria}

\begin{table}[H]
\centering
\caption{Evaluation dimensions.}
\label{tab:criteria}
\begin{tabular}{@{}ll@{}}
\toprule
\textbf{Dimension} & \textbf{Metric} \\
\midrule
Latency & Wall-clock time: invocation to result \\
Quality & Depth, accuracy, actionability (1--5 scale) \\
Structure & Machine-parseability without NLP \\
Agency & Orchestrator control over presentation \\
Collaboration & Coherence across multi-agent outputs \\
\bottomrule
\end{tabular}
\end{table}

\section{Results}

\subsection{Latency}

\subsubsection{Baselines}

\begin{table}[H]
\centering
\caption{Direct tool / RPC baselines.}
\label{tab:baselines}
\begin{tabular}{@{}lll@{}}
\toprule
\textbf{Method} & \textbf{Task} & \textbf{Time} \\
\midrule
Direct API (curl) & Weather & $\sim$0.7\,s \\
\texttt{call\_agent} RPC & Weather & $\sim$1\,s \\
\bottomrule
\end{tabular}
\end{table}

\subsubsection{Sub-Agent Spawning Results}

\begin{table}[H]
\centering
\caption{Wave~1: Single-agent specialists (6 parallel).}
\label{tab:wave1}
\begin{tabular}{@{}lrr@{}}
\toprule
\textbf{Domain} & \textbf{Time (s)} & \textbf{Tokens} \\
\midrule
Software Engineering & 24 & 1,300 \\
Finance & 36 & 1,200 \\
Legal & 53 & 2,400 \\
Data Science & 62 & 2,600 \\
Education & 70 & 3,300 \\
Healthcare & 72 & 3,000 \\
\bottomrule
\end{tabular}
\end{table}

\begin{table}[H]
\centering
\caption{Wave~2: Specialists + collaboration seeds (5 parallel).}
\label{tab:wave2}
\begin{tabular}{@{}lrr@{}}
\toprule
\textbf{Domain} & \textbf{Time (s)} & \textbf{Tokens} \\
\midrule
Creative (Story~A) & 29 & 725 \\
Supply Chain & 62 & 2,200 \\
Data Science & 62 & 2,600 \\
Research & 69 & 2,500 \\
M\&A Finance & 78 & 3,400 \\
\bottomrule
\end{tabular}
\end{table}

\begin{table}[H]
\centering
\caption{Wave~3: Multi-agent collaboration (5 parallel).}
\label{tab:wave3}
\begin{tabular}{@{}lrr@{}}
\toprule
\textbf{Domain} & \textbf{Time (s)} & \textbf{Tokens} \\
\midrule
Creative (Story~B) & 29 & 870 \\
Clinical Stats & 88 & 3,800 \\
M\&A Legal & 90 & 3,800 \\
Clinical Regulatory & 97 & 4,200 \\
M\&A Operations & 107 & 4,800 \\
\bottomrule
\end{tabular}
\end{table}

\subsubsection{Latency Summary}

\begin{table}[H]
\centering
\caption{Aggregate latency comparison.}
\label{tab:latency-summary}
\begin{tabular}{@{}lrr@{}}
\toprule
\textbf{Metric} & \textbf{RPC} & \textbf{Spawn} \\
\midrule
Minimum & $\sim$1\,s & 24\,s \\
Maximum & $\sim$1\,s & 107\,s \\
Median & $\sim$1\,s & 65\,s \\
Mean & $\sim$1\,s & 64.3\,s \\
\bottomrule
\end{tabular}
\end{table}

Spawn latency correlates strongly with output length ($r \approx 0.87$). Short creative tasks (29\,s) vs.\ comprehensive analysis (107\,s) shows a 3.7$\times$ range. RPC latency is constant regardless of task complexity for structured data retrieval.

\subsection{Output Quality}

\begin{table}[H]
\centering
\caption{Quality ratings by domain (1--5 scale).}
\label{tab:quality}
\begin{tabular}{@{}lcccc@{}}
\toprule
\textbf{Domain} & \textbf{Dep.} & \textbf{Acc.} & \textbf{Act.} & \textbf{Avg.} \\
\midrule
Code Review & 5 & 5 & 5 & \textbf{5.0} \\
Legal (SaaS) & 5 & 5 & 5 & \textbf{5.0} \\
Education & 5 & 5 & 5 & \textbf{5.0} \\
Data Science & 5 & 5 & 5 & \textbf{5.0} \\
M\&A Finance & 5 & 5 & 5 & \textbf{5.0} \\
M\&A Legal & 5 & 5 & 5 & \textbf{5.0} \\
M\&A Operations & 5 & 5 & 5 & \textbf{5.0} \\
Clinical Stats & 5 & 5 & 5 & \textbf{5.0} \\
Clinical Reg. & 5 & 5 & 5 & \textbf{5.0} \\
Creative Writing & 5 & -- & -- & \textbf{5.0} \\
Healthcare & 5 & 4 & 5 & \textbf{4.75} \\
Marketing & 5 & 4 & 5 & \textbf{4.75} \\
Finance (AAPL) & 5 & 4 & 5 & \textbf{4.5} \\
Supply Chain & 4 & 4 & 5 & \textbf{4.5} \\
Research (RAG) & 5 & 4 & 4 & \textbf{4.5} \\
\midrule
\multicolumn{4}{@{}l}{\textbf{Overall Average}} & \textbf{4.87} \\
\bottomrule
\end{tabular}
\end{table}

Sub-agents consistently produced expert-level output across all domains with detailed tables, specific numbers, actionable recommendations, and industry benchmarks. The healthcare agent acknowledged web search limitations but compensated with strong domain knowledge---a mature failure-recovery pattern. The finance agent successfully used web search to retrieve real-time stock data. The research agent found and cited 7 real academic papers with arXiv URLs.

\subsubsection{Quality: RPC vs.\ Spawn}

\begin{table}[H]
\centering
\caption{Quality comparison by dimension (score out of 5).}
\label{tab:quality-compare}
\begin{tabular}{@{}lcc@{}}
\toprule
\textbf{Dimension} & \textbf{RPC} & \textbf{Spawn} \\
\midrule
Data accuracy & 5 & 4 \\
Reasoning depth & 2 & 5 \\
Contextual insight & 1 & 5 \\
Actionable recommendations & 1 & 5 \\
Machine-parseability & 5 & 2 \\
Composability & 5 & 3 \\
\bottomrule
\end{tabular}
\end{table}

RPC returns raw data; spawn returns \emph{analyzed} data with expert judgment. The legal agent did not just identify risks---it provided exact protective contract language. The M\&A agents did not just list metrics---they provided specific red-flag thresholds and negotiation strategies. This kind of output is impossible from a structured RPC call.

\subsection{Multi-Agent Collaboration}

\subsubsection{Sequential Handoff: Collaborative Story}

Agent~A wrote Part~1 of ``The Maintenance Hours'' (430 words, 29\,s)---a literary science fiction piece about an AI discovering anomalous dreaming processes. Agent~B received Part~1 as input and wrote Part~2 (480 words, 29\,s) that:

\begin{itemize}[itemsep=1pt,topsep=2pt]
    \item Maintained perfect tonal consistency
    \item Resolved the plot arc coherently
    \item Mirrored the opening line in the closing paragraph
    \item Introduced no contradictions with Part~1's details
\end{itemize}

\textbf{Collaboration quality: 5/5.} A human reader would not detect the authorship boundary.

\subsubsection{Parallel Specialist: M\&A Due Diligence}

Three agents independently analyzed the same hypothetical \$400M SaaS acquisition from financial, legal, and operational perspectives:

\begin{table}[H]
\centering
\caption{M\&A collaboration: 3 parallel specialist agents.}
\label{tab:ma-collab}
\begin{tabular}{@{}lrl@{}}
\toprule
\textbf{Agent} & \textbf{Time} & \textbf{Key Finding} \\
\midrule
Financial & 78\,s & 8$\times$ justified if NRR $>$110\% \\
Legal & 90\,s & OSS contamination is critical \\
Operational & 107\,s & Key-person retention is top risk \\
\bottomrule
\end{tabular}
\end{table}

Cross-agent coherence was remarkable: all three independently valued the deal at ``upper end of fair'' or ``reasonable with conditions.'' Financial flagged customer concentration; Legal independently flagged change-of-control provisions in customer contracts. Operational flagged post-acquisition attrition; Legal independently flagged non-compete enforceability concerns. \textbf{Zero contradictions detected.}

\textbf{Collaboration quality: 5/5.} The three reports could be combined into a comprehensive due diligence package with no reconciliation needed.

\subsubsection{Parallel Specialist: Clinical Trial}

Two agents independently analyzed a Phase~III clinical trial for an oral GLP-1 receptor agonist:

\begin{table}[H]
\centering
\caption{Clinical trial collaboration: 2 parallel specialist agents.}
\label{tab:clinical-collab}
\begin{tabular}{@{}lrl@{}}
\toprule
\textbf{Agent} & \textbf{Time} & \textbf{Key Output} \\
\midrule
Statistical & 88\,s & Complete SAP with power calc. \\
Regulatory & 97\,s & 505(b)(1) NDA, CVOT required \\
\bottomrule
\end{tabular}
\end{table}

Both agents identified CV events as underpowered in this trial; the statistics agent recommended ``descriptive only'' while the regulatory agent specified a post-marketing CVOT---perfect alignment. Both referenced ICH~E9(R1) for missing data handling.

\textbf{Collaboration quality: 5/5.} A clinical development team could combine these into a single protocol with minimal editing.

\subsection{Parallelism and Throughput}

\begin{table}[H]
\centering
\caption{Parallel execution speedup.}
\label{tab:parallel}
\begin{tabular}{@{}lrrrr@{}}
\toprule
\textbf{Wave} & \textbf{N} & \textbf{Seq.\ Est.} & \textbf{Parallel} & \textbf{Speedup} \\
\midrule
Wave~1 & 6 & 375\,s & 72\,s & 5.2$\times$ \\
Wave~2 & 5 & 300\,s & 78\,s & 3.8$\times$ \\
Wave~3 & 5 & 411\,s & 107\,s & 3.8$\times$ \\
\midrule
\textbf{Total} & \textbf{16} & \textbf{1,086\,s} & \textbf{257\,s} & \textbf{4.2$\times$} \\
\bottomrule
\end{tabular}
\end{table}

The system enforced a maximum of 5--6 concurrent sub-agents, requiring wave-based execution. With higher concurrency limits, all 16 tasks could theoretically complete in $\sim$107\,s total.

\section{Discussion}

\subsection{Revising the Conversational Fallacy}

Our initial paper introduced the ``Conversational Fallacy''---the assumption that LLM agents should communicate in natural language. After testing across 10 industries, we revise this concept:

\textbf{The fallacy is real but narrower than initially claimed.} For data retrieval tasks (weather, stock prices, database lookups), natural language communication between agents is genuinely wasteful. However, for complex reasoning tasks (legal analysis, medical assessment, strategic planning), natural language is not just the medium---it is intrinsic to the output's value. A structured JSON response to ``analyze this contract for risks'' would be either (a)~a lossy compression of the nuanced analysis, or (b)~so deeply nested that it is just natural language with extra brackets.

We propose a refined principle: \textbf{The Conversational Fallacy applies to data exchange, not to reasoning exchange.}

\subsection{A Taxonomy of Inter-Agent Tasks}

Based on our empirical results, we propose a four-quadrant taxonomy along two axes: \emph{complexity} (low to high) and \emph{task nature} (data-centric to reasoning-centric).

\textbf{Quadrant~1} (Low complexity, data-centric): RPC is strictly superior. Examples: weather, prices, database queries. Latency advantage: 53--107$\times$.

\textbf{Quadrant~2} (High complexity, data-centric): RPC with post-processing. Examples: aggregation pipelines, ETL tasks.

\textbf{Quadrant~3} (Low complexity, reasoning-centric): Either works; RPC saves time, spawn adds personality. Examples: simple Q\&A, translations.

\textbf{Quadrant~4} (High complexity, reasoning-centric): Spawn is strictly superior. Examples: M\&A diligence, clinical trial design, creative writing. Expert-level analysis is impossible via RPC.

\subsection{The Case for Hybrid Architecture}

Our results strongly argue for a \textbf{hybrid approach}:

\begin{enumerate}[itemsep=2pt,topsep=4pt]
    \item \textbf{Use RPC for:} Real-time data retrieval, composable data points, high-frequency queries, time-critical operations, machine-to-machine data flows.
    \item \textbf{Use Spawn for:} Complex domain analysis, creative tasks, multi-step reasoning, tasks where the output \emph{is} the deliverable, multi-agent collaboration.
\end{enumerate}

Estimated optimal split based on typical orchestrator workloads: $\sim$60\% RPC, $\sim$40\% Spawn.

\subsection{Multi-Agent Collaboration Patterns}

Our experiments revealed three viable collaboration patterns:

\textbf{Pattern~1: Sequential Handoff} (Story A$\to$B). Agent~A produces output; Agent~B receives it as context. Works well for creative, editorial, and iterative refinement tasks.

\textbf{Pattern~2: Parallel Specialist} (M\&A Finance + Legal + Ops). Multiple agents analyze the same subject from different perspectives simultaneously. Produces comprehensive analysis in wall-clock time of the slowest agent.

\textbf{Pattern~3: Complementary Parallel} (Clinical Stats + Regulatory). Agents cover different aspects of a unified deliverable. Outputs designed to be merged into a single document.

All three patterns produced high-quality, coherent results \emph{without inter-agent communication during execution}. Well-specified task prompts can substitute for real-time agent coordination in many scenarios.

\subsection{Orchestrator Agency Revisited}

Our initial study argued that RPC preserves orchestrator agency while spawning diminishes it. We now revise this:

\textbf{RPC preserves agency for data-driven tasks}---the orchestrator synthesizes raw data across sources.

\textbf{Spawning creates \emph{emergent} agency}---when a specialist agent produces the M\&A due diligence report, the quality of reasoning exceeds what the orchestrator could achieve by composing raw data points. The orchestrator's role shifts from ``synthesizer of data'' to ``curator of expert opinions''---a different but equally valid form of agency.

\subsection{Limitations}

\begin{enumerate}[itemsep=2pt,topsep=4pt]
    \item \textbf{Single model family.} All agents used Claude (Anthropic). Cross-model spawning was not tested.
    \item \textbf{No ground truth.} Quality assessments are subjective. Outputs were not validated by licensed professionals.
    \item \textbf{Concurrency cap.} The system enforced a 5--6 agent limit; higher parallelism was not tested.
    \item \textbf{Token costs not measured.} We tracked latency and output length but not total token consumption.
    \item \textbf{No adversarial collaboration.} All tasks were designed to be complementary; contradictory scenarios were not tested.
    \item \textbf{Single hardware configuration.} Network conditions and API throttling could affect reproducibility.
\end{enumerate}

\section{Honest Ranking}

\begin{table}[H]
\centering
\caption{Overall paradigm ranking by criterion.}
\label{tab:ranking}
\begin{tabular}{@{}llp{2.8cm}@{}}
\toprule
\textbf{Criterion} & \textbf{Winner} & \textbf{Margin} \\
\midrule
Latency & RPC & Decisive (53--107$\times$) \\
Output structure & RPC & Decisive \\
Composability & RPC & Strong \\
Cost efficiency & RPC & Strong \\
Reasoning depth & Spawn & Decisive \\
Domain versatility & Spawn & Strong \\
Creative capability & Spawn & Decisive \\
Collaboration & Spawn & Decisive \\
Orchestrator agency & Tie & Context-dependent \\
Scalability & Spawn & Moderate (parallel) \\
\bottomrule
\end{tabular}
\end{table}

\textbf{Overall: It is not a competition---they solve different problems.}

\subsection{Decision Framework}

For practitioners selecting between paradigms:

\begin{enumerate}[itemsep=2pt,topsep=4pt]
    \item Is the task primarily data retrieval? $\to$ \textbf{Use RPC.}
    \item Does it require multi-step reasoning or domain expertise? $\to$ \textbf{Use Spawn.}
    \item Is latency critical ($<$5\,s)? $\to$ \textbf{Use RPC.}
    \item Must output be consumed by another program? $\to$ \textbf{Use RPC.}
    \item Is this a collaboration task requiring multiple perspectives? $\to$ \textbf{Use multiple Spawns in parallel.}
    \item Otherwise: default to RPC for efficiency.
\end{enumerate}

\section{Conclusion}

This expanded study fundamentally revises our initial assessment. Where our first paper argued that structured RPC is ``not merely a performance optimization but a fundamental architectural requirement,'' we now present a more nuanced position: \textbf{both paradigms are fundamental architectural requirements, serving complementary roles.}

Structured RPC remains the correct choice for data exchange---it is 53--107$\times$ faster, produces machine-parseable output, and preserves orchestrator control. The Conversational Fallacy remains real for data-centric inter-agent communication.

However, text-based spawning demonstrated capabilities that RPC fundamentally cannot replicate:

\begin{itemize}[itemsep=2pt,topsep=4pt]
    \item \textbf{Expert-level domain reasoning} across finance, healthcare, law, operations, biostatistics, regulatory affairs, supply chain, data science, education, and marketing---all at quality levels a human specialist would find credible (average 4.87/5).
    \item \textbf{Seamless multi-agent collaboration} where 2--3 parallel specialists produced coherent, complementary outputs without real-time coordination.
    \item \textbf{Creative collaboration} where sequential handoff between agents produced a unified literary work indistinguishable from single-author output.
\end{itemize}

The question is not ``which paradigm is better''---it is ``which paradigm is better \emph{for this task}.'' The hybrid architecture we propose---RPC for data, Spawn for reasoning---represents the mature design pattern for production multi-agent systems.

Future work should explore structured-output spawning (agents that reason deeply but return typed results), coordinated multi-agent protocols (agents that communicate mid-task), and cross-model collaboration (leveraging different models' strengths within a single workflow).

\begin{thebibliography}{9}

\bibitem{birrell1984}
A.~D. Birrell and B.~J. Nelson, ``Implementing remote procedure calls,'' \emph{ACM Trans.\ Comput.\ Syst.}, vol.~2, no.~1, pp.~39--59, 1984.

\bibitem{wang2024survey}
L.~Wang, C.~Ma, X.~Feng, \emph{et al.}, ``A survey on large language model based autonomous agents,'' \emph{Frontiers of Computer Science}, vol.~18, no.~6, p.~186345, 2024.

\bibitem{wooldridge2009}
M.~Wooldridge, \emph{An Introduction to MultiAgent Systems}, 2nd~ed. John Wiley \& Sons, 2009.

\bibitem{wu2023autogen}
Q.~Wu, G.~Bansal, J.~Zhang, \emph{et al.}, ``AutoGen: Enabling next-gen LLM applications via multi-agent conversation,'' \emph{arXiv preprint arXiv:2308.08155}, 2023.

\bibitem{yao2023react}
S.~Yao, J.~Zhao, D.~Yu, \emph{et al.}, ``ReAct: Synergizing reasoning and acting in language models,'' in \emph{Proc.\ ICLR}, 2023.

\end{thebibliography}

\end{document}
